\scanpage{23}

\[
\mathcal A_Q\simeq\operatorname{\unsure{Duaux|Deux}}_{\underline O}(\Lambda).
\]

L'hom

\[
u'_Q:\mathcal A_Q\mathbin{\unsure{-|\longrightarrow}}\underline{\operatorname{End}}(E)
\]

est \deleted{défini}\inserted{déterminé} par la condition \deleted{que} de s'annuler sur $\Lambda$, i.e. c'est $(a,\lambda)\longmapsto a$. \deleted{\unsure{Elle}} Il n'a rien d'injectif.

2.6. Quand on se donne une base de $E$ et de $\Delta$, alors $Q$ est donnée par

\[
Q(\lambda,\mu)=\alpha\lambda^{\unsure{1|2}}+\beta\lambda\mu+\gamma\mu^2
\qquad\alpha,\beta,\gamma\in\Gamma\underline O,
\]

et on \deleted{peut prendre} a une \deleted{\unsure{syst.|base}}\unsure{base} $(1,u)$ de $\mathcal A_Q$, avec

\[
\left|\begin{aligned}
u^2-\beta u+\alpha\gamma&=0\qquad\text{i.e. }\deleted{\beta=\operatorname{Tr}u},\\
\operatorname{Tr}u&=\beta\\
Nu&=\alpha\gamma
\end{aligned}\right.
\]

L'image de $u$ dans $\operatorname{End}(E)\simeq\mathbb M_2(\underline O)$ est la matrice déjà envisagée

\[
u_Q=\begin{pmatrix}0&-\gamma\\\alpha&\beta\end{pmatrix}
\]

Pour achever de déterminer $\mathcal A_Q$ comme foncteur en $(E,\Delta,Q)$, il faut voir comment s'explicite l'isom de $A_{\alpha,\beta,\gamma}=A[u]/(u^2-\beta u+\alpha\gamma)$ avec $A_{\alpha',\beta',\gamma'}=A[u]/(u^2-\beta'u+\alpha'\gamma')$, quand on passe d'un syst. de bases de $(A,\Delta)$ à un autre, par des matrices $\begin{pmatrix}u_{11}&u_{12}\\u_{21}&u_{\unsure{12|22}}\end{pmatrix}$ et $v$ \unsure{..|\ldots}
% archive sheet number: 22
